% Generated draft: requires mathematical and source review.
\documentclass[11pt]{article}
\usepackage[margin=25mm]{geometry}
\usepackage{fontspec}
\setmainfont{Latin Modern Roman}
\usepackage{amsmath,amssymb,mathtools,unicode-math}
\setmathfont{Latin Modern Math}
\usepackage{enumitem,needspace,xurl,hyperref}
\setlength{\parindent}{0pt}
\setlength{\parskip}{5pt}
\setlength{\emergencystretch}{3em}
\newcommand{\N}{\mathbb{N}}
\newcommand{\Z}{\mathbb{Z}}
\newcommand{\Q}{\mathbb{Q}}
\newcommand{\R}{\mathbb{R}}
\newcommand{\C}{\mathbb{C}}
\newcommand{\F}{\mathbb{F}}
\newcommand{\A}{\mathbb{A}}
\newcommand{\PP}{\mathbb{P}}
\newcommand{\E}{\mathbb{E}}
\newcommand{\eps}{\varepsilon}
\newcommand{\dd}{\,\mathrm{d}}
\begin{document}
\section*{UnsolvedMath ID: 20002470}
\textbf{Source problem code:} catalogue entry
\subsection{Definitions and mathematical statement}
Write $\mathbb T=\R/\Z$. Mather's question asks: if a real-analytic twist map $f:\mathbb T\times\R\to\mathbb T\times\R$ has a foliation by invariant curves $\Gamma_c$ which are Aubry--Mather sets, must every $\Gamma_c$ be real analytic? The question concerns regularity of individual leaves, not necessarily analytic dependence of the entire foliation on $c$.

The aid uses the following standard but stronger explicit interpretation. Assume $f$ is an exact symplectic positive twist diffeomorphism, isotopic to the identity, with lift $F=(X,Y)$ satisfying $\partial_yX>0$ and
\[
 f^*(y\,\dd x)-y\,\dd x=\dd S
\]
for a globally defined function $S$ on the cylinder. Suppose the leaves are whole graphs
\[
 \Gamma_c=\{(x,g_c(x)):x\in\mathbb T\},\qquad
 \int_{\mathbb T}g_c(x)\,\dd x=c,
\]
forming a continuous foliation, and each graph is the entire $c$-Aubry set of the associated variational problem. Here the $c$-modified discrete action subtracts $c$ times the lifted angular displacement; the Aubry condition includes the appropriate chain-recurrence condition, rather than merely an arbitrary invariant graph of minimizing orbits. The original wording does not explicitly fix exactness, twist sign, foliation topology, or which Aubry/Mather convention is intended.

Let $\rho(c)$ be the rotation number of the circle homeomorphism induced by $f$ on $\Gamma_c$, computed using the chosen lift. The supplied aid reports:
\begin{enumerate}
\item If $\rho(c)=p/q\in\Q$ with $q>0$, every point of the whole Aubry leaf is periodic and
\[
 F^q(x,g_c(x))=(x+p,g_c(x)).
\]
The minimizing no-conjugate-points argument makes the fixed-displacement equation suitable for the analytic implicit-function theorem, so $g_c$ is real analytic.
\item The map $c\mapsto\rho(c)$ is continuous and nondecreasing, with every nontrivial plateau having rational value. Consequently parameters with rational rotation number are dense, and the corresponding analytic leaves form a dense rational skeleton of the foliation.
\end{enumerate}
These are reported partial results under the whole-Aubry-foliation interpretation. They do not prove analyticity of an irrational leaf: uniform real convergence of analytic graphs alone is insufficient. One sufficient additional estimate would be a common complex strip $|\operatorname{Im}z|<a$, $a>0$, on which a sequence of rational approximating graph functions has uniformly bounded holomorphic extensions. Together with convergence to $g_c$ on the real circle, normal-family compactness would then yield analyticity of that limit. Neither such uniform complex control for all irrational leaves nor a counterexample with degenerating control is supplied. No nonlinear foliation-conjugacy or global analytic-integrability conclusion is asserted.
\subsection{Short English statement}
Ask whether every invariant Aubry--Mather leaf of an analytic twist-map foliation is analytic. Under the aid's exact positive-twist, whole-Aubry interpretation, rational leaves give a dense analytic skeleton; the missing issue is analytic control at irrational limits.
\subsection{Sources}
\begin{itemize}
\item John Mather, Problem 7, American Institute of Mathematics workshop \emph{New connections between dynamical systems and PDE's}, \url{https://aimath.org/WWN/dynpde/dynpde.pdf}.
\item Local UnsolvedMath registry, permanent ID 20002470, source code \texttt{catalogue entry}, canonical location \texttt{aim-pdes-notes.json}, \texttt{notes[81]}, record dated 14 August 2026. The original literature field is empty.
\item Supplied research attempt 1: partial result, problem status at run open, rigor labelled proved, novelty confidence low. Its rational-leaf and density reduction is machine-generated research aid, not peer reviewed, requiring independent verification before citation; no separate authoritative reference for those refinements is supplied.
\end{itemize}
\end{document}
